\documentclass[11pt]{article}

\usepackage[utf8]{inputenc}
\usepackage[T1]{fontenc}
\usepackage{lmodern}
\usepackage{upquote}
\usepackage{xcolor}
\usepackage{listings}
\usepackage[margin=1in]{geometry}
\usepackage{amsmath,amssymb}
\usepackage{booktabs}
\usepackage{enumitem}
\usepackage[hidelinks]{hyperref}
\usepackage{url}

\definecolor{codegray}{gray}{0.45}
\lstset{basicstyle=\ttfamily\small,breaklines=true,columns=fullflexible,
        xleftmargin=4pt,xrightmargin=4pt,showstringspaces=false}
\setlength{\emergencystretch}{3em}

\newcommand{\R}{\mathbb{R}}
\newcommand{\N}{\mathbb{N}}
\newcommand{\valen}{\textsc{Valen}}

\title{\bfseries From Structural Invariants to a Red-Team Platform:\\ \Large
Attack-Graph Mathematics for Autonomous Offensive Operations}
\author{\valen\ Project}
\date{Working draft --- engineering and methods report}

\begin{document}
\maketitle

\begin{abstract}
\noindent We describe the extension of \valen\ --- a neuro-symbolic
structural-invariant verifier --- into a full red-team \emph{platform}: a
multi-user team server that coordinates autonomous offensive operations across
the kill chain (reconnaissance, Active Directory attack-graph analysis, command
\& control, credential recovery, phishing and exfiltration), under an explicit
authorization and human-in-the-loop gating model. The core contribution is a
layered attack-graph \emph{mathematics}: structural signals (spectral, persistent
and directed homology, discrete curvature), \emph{flow} (Menger min-cut
chokepoints, min-cost flow, articulation points), \emph{probability} (absorbing
Markov-chain hitting probabilities, solved exactly by $(I-Q)^{-1}R$),
\emph{formal synthesis} (Z3 Optimize bounded model checking that yields minimal-cost
ordered exploit plans with AND-preconditions), and a \emph{probabilistic
context-free grammar} for password guessing. We report the system as deployed
against OWASP crAPI (18/18 challenges solved autonomously, including three LLM
prompt-injection challenges) and the Active Directory toy corpus, and we state
precisely what each mathematical layer answers that a naive shortest-path or
reachability query cannot.
\end{abstract}

\section{Introduction}
Classic semantic analysis answers \emph{is this sink reachable by untrusted
data?}. Offensive operations ask harder, quantitative questions:
\begin{itemize}[itemsep=1pt]
  \item Which principals are \emph{obligatory} on every path to the target (the
        chokepoints)?
  \item What is the \emph{cheapest} ordered sequence of actions, given a model of
        stealth and effort, not just the fewest-hop one?
  \item What is the \emph{probability} that a given entry point reaches the
        crown-jewel target, when the attacker moves according to some policy?
  \item With a \emph{budget} of $k$ compromises, what are the $k$ cheapest
        independent routes?
  \item Which passwords should a spray try first, to maximize hits before lockout?
\end{itemize}
This paper reports the machinery added to \valen\ to answer these, and the
platform that operationalizes it.

\section{Platform architecture}
\valen\ is a single installable Python package (\texttt{pip install valen}) with
an optional Rust numeric core and a pure-Python fallback. Its server is a
multi-user \emph{team server} (users with roles \texttt{viewer}/\texttt{operator}/
\texttt{admin}, PBKDF2 sessions, per-route RBAC, Server-Sent Events for live
collaboration) that persists engagements, runs, and evidence artifacts in SQLite.

The red-team surface spans the kill chain:
\begin{itemize}[itemsep=1pt]
  \item \textbf{Recon / enumeration} --- nmap, masscan, gobuster, ffuf, nuclei
        wrappers with stealth profiles and CVE/KEV/EPSS intelligence.
  \item \textbf{Active Directory} --- BloodHound/SharpHound ingestion into a typed
        IR graph, tier-0 attack paths, Kerberoast/AS-REP selection, GPP
        \texttt{cpassword} decryption, and command builders.
  \item \textbf{C2 and payloads} --- Sliver integration (gRPC-first, shell
        fallback), msfvenom/handler/HTA planners, phishing (GoPhish) and
        exfiltration planners.
  \item \textbf{Hybrid agent} --- a deterministic kill-chain planner with an
        operator approval queue; intrusive actions require explicit
        \texttt{--authorize}.
  \item \textbf{Reporting} --- HTML / Markdown / JSON / SARIF with CVSS 3.1/4.0
        and CVE/KEV/EPSS enrichment.
\end{itemize}
Everything offensive is gated by a scope allow-list (loopback by default), an
execution allow-list (\texttt{--allow-exec}), and per-operator roles.

\section{Attack-graph mathematics}
The IR is a typed directed multigraph $G=(V,E)$ over principals. We layer four
families of methods over it.

\subsection{Structural signals}
The base layer reuses the existing kernels: the combinatorial Laplacian and its
Fiedler vector/$\lambda_2$, persistent $H_0/H_1$ homology, GLMY directed path
homology, Forman--Ollivier Ricci curvature, and betweenness/PageRank. To these we
add spectral centralities --- eigenvector, Katz, and HITS
(hubs/authorities) --- for targeting and pivot selection.

\subsection{Flow: chokepoints and disjoint plans}
\textbf{Menger min-cut.} Given entry set $S$ and target set $T$, we compute the
minimum vertex set whose removal separates them. We reduce to an $s$-$t$ max flow
on the vertex-split graph (each vertex $v$ becomes $v_{\mathrm{in}}\to
v_{\mathrm{out}}$ of capacity $1$; sources/targets unbounded) and run
Edmonds--Karp; the saturated split edges on the source side of the residual graph
are the \emph{obligatory chokepoints}. A flow of magnitude $\geq \infty$ means a
source is directly adjacent to a target (min-cut $0$).

\textbf{Min-cost flow.} With a per-edge cost model (stealth + effort; e.g.\
$\text{MemberOf}=1,\ \text{DCSync}=5$), we send $k$ units of flow at minimum cost
via successive shortest augmenting paths (Bellman--Ford on the residual graph).
Because internal vertices have capacity $1$, the result is a set of
\emph{vertex-disjoint} cheapest attack routes --- ``with $k$ compromises, the $k$
cheapest independent ways in''.

\textbf{Articulation points and bridges} (Tarjan, on the symmetrized graph) rank
the most disruptive principals and links: removing them fragments the network.
The directed min-cut is the \emph{defender/hardening} view; the undirected
articulation points are the \emph{attacker/disruption} view.

\subsection{Probability: hitting time}
Model the attack graph as an absorbing Markov chain: targets are absorbing, and
each transient node transitions uniformly over its out-neighbours. The hitting
probability $h(u)$ satisfies
\[
h(u) \;=\; \mathbb{1}[u \in T] \;+\; \mathbb{1}[u\notin T]\cdot
   \frac{1}{\deg^+(u)}\sum_{v : u\to v} h(v),
\]
computed either by value iteration or \emph{exactly} by Gaussian elimination on
$(I-Q)h = r$, where $Q$ is the transient-to-transient transition and $r$ the
one-step mass into targets. This yields a probabilistic risk ranking, distinct
from boolean reachability: in the toy corpus, an entry with three out-edges and a
single route to the target scores $\tfrac13$, while a direct entry scores $1$.

\subsection{Formal synthesis: minimal-cost exploit plans}
Boolean reachability (Z3) answers \emph{is it reachable?}. We upgrade to
\emph{synthesis}: model the graph as a transition system
\[
\mathrm{has}(s{+}1, n) \;\equiv\; \mathrm{has}(s,n)
   \;\vee\; \bigvee_{e:\, m\to n} \bigl(\mathrm{taken}_e \wedge \mathrm{has}(s,m)\bigr),
\]
with entries compromised at step $0$, and ask \texttt{z3.Optimize} to minimize
$\sum_e \mathrm{taken}_e \cdot \mathrm{cost}(e)$ subject to reaching the target
within $k$ steps. Reconstructing the ordered action sequence gives the
technique-by-technique narrative. We extend this to \textbf{AND-preconditions}:
compound actions $\{\mathrm{sources}\to \mathrm{target}\}$ that fire only when
\emph{all} sources are compromised (e.g.\ DCSync needs Domain Admins \emph{and} a
session on a Domain Controller).

\subsection{Password guessing: PCFG}
We implement Weir's probabilistic context-free grammar: tokenize passwords into
$L/D/S$ runs; the base structure (e.g.\ $L_6D_4S_1$) is the non-terminal and
terminals are keyed by $(\text{class},\text{length})$ --- so no cross-length
mixing ($\texttt{summer}+123$ is rejected when only $L_6D_4$ was seen). The
probability of a guess is $P(\mathrm{structure})\cdot\prod P(\mathrm{terminal})$,
and guesses are emitted in strictly decreasing probability by a best-first
(Dijkstra) walk over the sorted per-slot terminals. Training corpus is a cracked
hashcat potfile; an empty corpus falls back to a seed list.

\section{Evaluation}
\subsection{Autonomous pentest against OWASP crAPI}
Against a fresh OWASP crAPI 1.1.5 lab, the autonomous agent solves \textbf{18/18}
documented challenges, including BOLA/BFLA, mass assignment, SSRF, NoSQL/SQL
injection, JWT forgery, and three LLM prompt-injection challenges (routed through
a local LiteLLM proxy). The BOLA detector reaches recall $1.0$ / precision $0.90$.

\subsection{Active Directory toy corpus}
On a synthetic SharpHound corpus the layers compose:
\begin{itemize}[itemsep=1pt]
  \item \emph{Chokepoints}: with entry \texttt{BOB}, min-cut $1$ at the
        intermediate group \texttt{IT}; with entry \texttt{ALICE} (directly
        GenericAll on Domain Admins), min-cut $0$.
  \item \emph{Hitting}: $\texttt{ALICE}=1.0$, $\texttt{BOB}=\tfrac13$ (three
        out-edges, one to the target).
  \item \emph{Synthesis}: \texttt{ALICE} $\xrightarrow{\text{GenericAll/T1098}}$
        Domain Admins (cost $3.0$), and \texttt{ALICE}
        $\xrightarrow{\text{MemberOf}}\,$ \texttt{IT}
        $\xrightarrow{\text{DCSync/T1003.006}}$ domain (cost $6.0$); \texttt{BOB}
        is certified \emph{unreachable} to Domain Admins.
\end{itemize}

\section{Limitations and ethics}
The platform is an \emph{authorized-engagement} tool: scope is loopback by
default, execution is behind \texttt{--allow-exec}, intrusive actions require
operator approval, and every module carries a "for authorized engagements only"
notice. The pure-Python numerical fallback is an honest approximation of the Rust
core (persistent homology reduced to Betti numbers, GLMY to a cycle-rank
estimate); set \texttt{VALEN\_CORE\_BIN} to use the compiled core. The PCFG prior
is only as good as its training corpus; production use should feed it a real
cracked corpus.

\section{Conclusion}
We turned a structural-invariant verifier into a red-team platform whose
differentiator is \emph{quantitative attack-graph mathematics}: chokepoints,
least-cost and disjoint plans, exact hitting probabilities, formal minimal-cost
exploit synthesis with AND-preconditions, and a PCFG password prior. Each layer
answers a question that reachability and shortest paths cannot, and the whole is
operationalized behind an explicit authorization model. All code is installable,
tested (383 tests), and documented.

\end{document}
